Isogeometric analysis employs unified spline basis functions to describe CAD geometry and analysis fields, thereby providing an integrated framework for geometric modeling, numerical analysis and sensitivity evaluation in size, shape and topology optimization. For complex shells, curved cellular structures and additive-manufactured lightweight components, conventional workflows based on repeated geometry reconstruction, remeshing and optimization updates often suffer from frequent model transfer, limited boundary quality and inconsistent design-variable semantics. Isogeometric analysis can organize analysis, optimization updates and geometry feedback within a unified geometric representation, which provides a technical basis for high-fidelity structural optimization. This review focuses on isogeometric analysis for intelligent structural optimization. It first introduces the basic B-spline/NURBS formulation and isogeometric discretization, and then summarizes classical models including SIMP material interpolation, adjoint sensitivity analysis, filtering regularization, level-set evolution and free-form-deformation parameterization. Representative routes for isogeometric topology and shape optimization are reviewed from density-based methods, level-set methods, moving morphable components and multi-patch shell optimization. Recent trends in metamaterials, lattice structures, additive manufacturing, intelligent generation and open-source workflows are further discussed. Existing studies indicate that complex multi-patch geometries, local refinement, manufacturability constraints, geometry feedback and the coupling between data-driven methods and physics-based models remain key issues for engineering applications.
